\scanpage{68}
\mspage{33}
\uline{Proposition} The divisor at infinity \inserted{$M^\infty_{0,I}$} of $\widehat M_{0,I}$ is a \inserted{\unsure{split|split}} standard curve over $\operatorname{Spec}\mathbb Z$, whose graph is isomorphic to the \deleted{that of the} $1$-skeleton of the skew dodecahedron defined by $I$ (equipped with a suitable orientation).
\margin{\uline{NB} $\widehat M_{0,I}\longrightarrow$ a genuine projective and smooth scheme \unsure{and all|and so much} --- the groupoid it represents is indeed rigid. It \unsure{holds} for all the $M_{0,I}\simeq M^!_{0,\nu}$ whatever $\nu\in\mathbb N\ldots$}

Of course, the genera of the irreducible components are zero, There is no \inserted{need to put} any ``marked points'' on \deleted{\unsure{$M_{0,I}$}} the \inserted{$10$} irreducible components of $M^\infty_{0,I}$, which is in any case an MD-curve \deleted{\unsure{stable}} with maximal MD-graph (all genera zero, all vertices of order $3$, no marked feet).

The isomorphism of the preceding proposition takes into account the actions of $\mathfrak S_I^+$ \inserted{$\simeq\operatorname{Aut}\text{ of the skew dodecahedron}$} on $M^\infty_{0,I}$, but not those of $\mathfrak S_I$ \deleted{\unsure{all}} in its entirety --- one would have to understand it from the point of view of the geometry of the \inserted{skew} icosahedron (or of the \inserted{skew} dodecahedron\ldots) --- this must be related of course to the combinatorial skew \uline{bi-icosahedron}\unsure{)|]}

When one passes to the modular multiplicity
\margin{67}
